\begin{abstract}
Este trabajo aborda el problema de clasificación automática de imágenes de flores utilizando redes neuronales profundas. Se utiliza el conjunto de datos de la competencia \textit{Petals to the Metal -- Flower Classification on TPU} de Kaggle, compuesto por 23,847 imágenes RGB distribuidas entre 104 clases distintas de flores. El conjunto presenta variaciones importantes en tamaño, color, iluminación, orientación y composición de las imágenes, además de un desbalance marcado: en entrenamiento, la clase más pequeña tiene 18 imágenes y la más grande 782, con una mediana de 88.

La tarea consiste en desarrollar un modelo capaz de identificar correctamente la clase de una flor a partir de una imagen. Debido a la similitud visual entre algunas categorías y a la variabilidad dentro de una misma clase, se proponen dos arquitecturas concretas: una red neuronal convolucional (CNN) propia, entrenada desde cero como línea base, y una red preentrenada en ImageNet (EfficientNetB0, con ResNet50 como alternativa) adaptada mediante aprendizaje por transferencia, que es la opción más adecuada para el número reducido de imágenes por clase. Los métodos tradicionales (KNN, bosques aleatorios y SVM) se analizan como referencia y se explica por qué solo son viables sobre características extraídas y no sobre los 150,528 valores de cada imagen.

Para evaluar el desempeño se utiliza la métrica F1 macro de Kaggle, ya que asigna la misma importancia a cada una de las 104 clases independientemente de la cantidad de imágenes disponibles. De esta forma, el proyecto busca una estrategia de clasificación que logre un buen desempeño general sin favorecer únicamente a las clases con mayor cantidad de ejemplos.
\end{abstract}
